\documentclass[12pt,a4paper]{article}
\title{分析化学学习笔记}
\author{李明翰}
% =========================
% 前处理（建议使用 XeLaTeX 编译）
% =========================
\usepackage[UTF8]{ctex}
\usepackage{geometry}
\usepackage{amsmath,amssymb}
\usepackage{setspace}
\usepackage{chemfig}
\usepackage{indentfirst}
\usepackage{siunitx}
\usepackage{tikz}
\usetikzlibrary{arrows.meta,positioning}
\geometry{
  left=2.5cm,
  right=2.5cm,
  top=2.5cm,
  bottom=2.5cm
}

\setlength{\parindent}{2em}
\linespread{1.3}

\begin{document}
\maketitle
\newpage
\section{绪论}
\subsection{为什么要学分析化学}
分析化学是发展和应用各种方法、仪器和策略以获得有关物质在空间和时间方面组成和性质的科学。

分析化学是化学、生物、制药、医学的基础学科、工具学科。是有力地生活助手。培养科学的思维方法和严谨求实的工作习惯。
\subsection{分析方法的分类}
分析方法的分类：
\begin{itemize}
\item 按分析任务分类：定性分析、定量分析结构分析
\item 按分析对象分析：无机分析、有机分析、核分析
\item 按试样用量分类：微量分析、常量分析、半微量分析、超微量分析
\item 按待测成分含量分类：痕量分析、常量分析，微量分析
\item 按方法原理分类：仪器分析方法、化学分析方法
\item 按应用领域命名的方法：环境分析方法、食品分析方法、药物分析方法等
\item 特殊命名的方法：滴定分析方法、光谱分析方法、色谱分析方法、电化学分析方法等
\end{itemize}
\subsection{基本分析过程}
$${\text{取样}}\rightarrow{\text{试样的分解}}\rightarrow{\text{消除干扰}}\rightarrow{\text{分析测定}}\rightarrow{\text{分析结果}}$$
\subsubsection{试样的采集}
试样的采集是分析化学的第一步，直接关系到分析结果的准确性和可靠性。试样采集的原则包括代表性、完整性、无污染和安全性。气体样品和液体样品的均匀性要优于固体样品，根据试样的状态和性质，采集方法可以分为以下几类：
\subsubsection*{气体试样的采集}
{\bf{分析对象：}}大气、压缩气体、气体反应物、工业废气、汽车尾气等。\\
{\bf{大气监测项目中的常见气体：}}二氧化硫、氮氧化物、一氧化碳、臭氧、挥发性有机物等。\\
{\bf{采集方法：}}用泵将气体冲入取样容器密封保存；用吸附剂吸附气体成分后进行分析；用化学试剂吸收气体成分后进行分析。\\
{\bf{大气样品：}}通常选择距离地面50 - 180cm，与人体呼吸区相近的位置进行采集，用直接法或者浓缩法取样。\\
{\bf{烟道气、废气中某些有毒污染物的分析:}把气体样品采入空瓶中或者大型注射器中。}大型容器中的气体均匀性可能不好，可以分别取每个部分取样之后摇匀
\subsubsection*{液体试样的采集}
\bf{1. 分析对象与分析指标}

\begin{itemize}
    \item \textbf{分析对象：}饮用水、工业用水、污水、饮料、工业溶剂、天然水体（江、河、湖、海、雨水、地下水）。
    
    \item \textbf{主要目的：}测定水中杂质的种类和数量，判断水质优劣，并评价其是否满足某种用途的要求。
    
    \item \textbf{水质指标：}主要分为物理指标、化学指标和微生物指标。
    
    \item \textbf{物理指标：}温度、气味、色度、浑浊度、总固体量、电导率等。
    
    \item \textbf{化学指标：}化学需氧量（COD）、pH、硬度、碱度、氯化物、硫酸盐等无机物，溶解气体、有机物、有毒物质及放射性物质等。
    
    \item \textbf{微生物指标：}微生物种类和数量，常作为判断生物污染程度的指标。
\end{itemize}

\bf{2. 采样方式}

\begin{itemize}
    \item \textbf{小容器中的液体物料：}在搅拌下直接用瓶子或取样管取样。
    
    \item \textbf{大容器中不容易摇匀的物料：}应从不同位置和深度分别取样，再混合均匀。
    
    \item \textbf{水管中或有泵水井中的水样：}取样前应先将水龙头或泵打开放水约 $10\sim15\,\mathrm{min}$，然后再用干净瓶子收集水样。
    
    \item \textbf{池塘、湖泊、江河中的水样：}
    \begin{itemize}
        \item 先选择好采样地点；
        \item 用采样器在不同深度各取一份水样；
        \item 可用于流域、剖面组成分析；
        \item 也可按设定比例混合均匀，作为分析试样，以测定该水体的平均值。
    \end{itemize}
    
    \item \textbf{液体物料的采样器：}塑料容器（适于微量元素分析）或玻璃瓶（适于有机物分子分析）。
\end{itemize}

\bf{3. 保存与稳定化处理}

\begin{itemize}
    \item 对于组成易发生变化的水体样品，如果不能在规定时间内立即完成分析测定，则必须采取适当的保护措施，以防止或减少存放期间试样组成的变化。
    
    \item \textbf{常用保护措施：}控制溶液的 pH、加入稳定剂、冷藏或冷冻、避光和密封等。
    
    \item 样品保存期限的长短，与待测物的稳定程度及保存方法密切相关。
    
    \item \textbf{常见保存方法举例：}
    \begin{itemize}
        \item 加入 HgCl$_2$：抑制细菌生长；适用于多种形态的氮、多种形态的磷及有机氯农药分析。
        
        \item 加入 HNO$_3$，调至 $pH<2$：防止金属离子沉淀；适用于多种金属离子的测定。
        
        \item 加入 H$_2$SO$_4$，调至 $pH<2$：抑制细菌生长，并与有机碱形成盐类；适用于有机水样（如 COD、油和油脂、有机碳）以及氨、胺类的测定。
        
        \item 加入 NaOH：与挥发性酸性化合物形成盐类；适用于氰化物、有机酸类的测定。
        
        \item 冷冻：可抑制细菌生长、减缓化学反应速率；适用于酸碱度、有机物、BOD、色、味、有机磷、有机氯、有机硫等项目。
    \end{itemize}
\end{itemize}
\subsubsection*{固体试样的采集}
\begin{itemize}
    \item \textbf{种类繁多：}如岩石、矿物、土壤、动植物、工业产品等。
    
    \item \textbf{均匀性较差：}采样时必须重点考虑代表性问题。
    
    \item \textbf{多数固体样品水溶性差：}通常需要采用必要的化学处理步骤后再进行分析。
\end{itemize}
\subsubsection*{生物试样的采集}
\begin{itemize}
    \item \textbf{样品特点：}组分数量多、浓度低、化学性质相似、易失活。
    
    \item \textbf{采集特点：}生物试样的组成随部位和时间不同而有较大差异，因此采集时要注意群体代表性、时效性和部位典型性，并防止生物污染。
    
    \item \textbf{新鲜样品的处理：}
    \begin{itemize}
        \item 对于需新鲜分析的动、植物试样，应立即处理并分析；
        \item 当天不能分析的新鲜试样，应暂时保存于冰箱中。
    \end{itemize}
    
    \item \textbf{宜采用新鲜试样分析的组分：}酚、亚硝酸、有机农药、维生素、蛋白质（特别是酶）、抗原、抗体、核酸、氨基酸等易转化、易降解或不稳定的组分。
    
    \item \textbf{需干样分析的植物试样：}
    \begin{itemize}
        \item 采集后要用清洁水洗净；
        \item 立即晾干或烘干、粉碎；
        \item 根据分析方法要求，过 $40\sim100$ 号筛；
        \item 混匀备用；
        \item 一般应保证试样经加工处理后保留约 $1\,\mathrm{kg}$ 的干重试样。
    \end{itemize}
    
    \item \textbf{动物试样：}如尿液、血液、唾液、胃液、乳液、粪便、指甲、骨、肉等，采集后应根据分析要求进行适当处理。
    
    \item \textbf{毛发和指甲样品的洗涤：}采集后先用中性洗涤剂清洗，再用纯水冲洗，之后依次用丙酮、乙醚、酒精或 EDTA 洗涤。
\end{itemize}
\subsubsection{试样的分解}
\subsubsection*{1. 基本概念}

湿法分析中，溶液具有较高的均匀性，目前大多数分析测定都在溶液中进行。若试样本身不是溶液状态，则需要通过适当的方法将其转化为溶液状态，这一过程称为\textbf{试样的分解}。

\subsubsection*{2. 试样分解的原则}

\begin{enumerate}
    \item 试样分解必须完全，处理后的溶液中不得残留原试样的细屑或粉末。
    \item 在分解过程中不能引入待测组分和干扰物质，也不能使待测组分损失。
    \item 应根据试样的组成与性质、待测组分的性质以及分析目的，选择合适的分解方法。
\end{enumerate}

\noindent 常用的分解方法有：溶解法、熔融法、半熔法、干式灰化法、湿式灰化法和微波辅助消解法等。

\subsubsection*{3. 常用分解方法}

\paragraph{（1）溶解法}
溶解法是指采用适当的溶剂，将试样直接分解并制成溶液的方法。该方法操作较简单、速度较快，主要用于无机物的分解。常用的无机溶剂有水、酸和碱等。

许多碱金属盐类、硝酸盐、铵盐以及大多数碱土金属盐和氯化物等都易溶于水。对于不溶于水的无机物，常采用酸、碱或混合酸作溶剂进行溶解。若常温下不易溶解，可适当加热；若稀酸、稀碱中仍不易溶解，则可使用浓酸、浓碱进行溶解。

\paragraph{（2）干式灰化法}
干式灰化法适用于有机物和生物试样的分解，可用于测定金属元素以及硫、磷、卤素等元素的含量。

其优点是一般不加或只加入少量试剂，因此可减少外来杂质的引入，方法也较为简便；缺点是某些元素可能挥发损失，且样品还可能与坩埚发生作用。

\begin{itemize}
    \item \textbf{马弗炉燃烧法：}将固体试样置于马弗炉（高温炉）中，在 \(400\sim700\,^\circ\mathrm{C}\) 下加热燃烧，由空气中的氧起氧化作用。燃烧后留下无机残渣。残渣通常再用少量盐酸或热浓硝酸浸取，然后定量转移到玻璃容器中。为提高灰化效率，有时可根据需要加入少量氧化性物质作为助剂，如硝酸镁等。对于液态或泥状的动植物组织，灰化前通常应先经蒸汽浴或轻度加热使其干燥。使用马弗炉时，应从室温开始逐渐升温，以防试样着火或起泡。

    \item \textbf{氧瓶燃烧法：}将试样包在定量滤纸内，用铂丝夹夹牢，放入充满氧气的锥形烧瓶中燃烧，燃烧产物用适当的吸收液吸收。试样中的卤素、硫、磷及某些金属元素分别转化为相应的离子或氧化物并进入吸收液中。与马弗炉燃烧法相比，氧瓶燃烧法分解更完全，燃烧产物的吸收液可直接用于元素分析，适用于少量试样的分解，且操作较为简便、快速。

    \item \textbf{燃烧法测定碳、氢：}在有机化合物中碳、氢元素的测定中常采用该法。方法是将试样置于铂舟中，加入适量金属氧化物作催化剂，然后在氧气流中充分燃烧，使碳定量转化为 \(\mathrm{CO_2}\)，氢定量转化为 \(\mathrm{H_2O}\)。燃烧生成的 \(\mathrm{CO_2}\) 和 \(\mathrm{H_2O}\) 分别用预先称重并装有吸收剂的吸收管吸收，根据吸收管质量的增加计算试样中碳和氢的含量。一般用烧碱石棉吸收 \(\mathrm{CO_2}\)，用高氯酸镁吸收 \(\mathrm{H_2O}\)。
\end{itemize}

\paragraph{（3）湿式灰化法}
湿式灰化法适用于有机物和生物试样的分解。该法将试样与具有氧化性的混合酸一起置于凯氏（Kjeldahl）烧瓶或烧杯中，在一定温度下进行消解，以破坏有机物，从而测定其中的金属元素以及硫、磷、卤素等元素的含量。

常用的氧化性酸有硝酸、高氯酸和硫酸等；常见混合酸体系有硝酸+\allowbreak 硫酸、硝酸+\allowbreak 高氯酸+\allowbreak 硫酸等。与干式灰化法相比，湿式灰化法的优点是速度较快，缺点是因加入试剂而可能引入杂质，因此应尽量使用高纯试剂。一般来说，干式灰化法约需 \(2\sim4\,\mathrm{h}\)，而湿式灰化法约需 \(30\sim60\,\mathrm{min}\)。

\paragraph{（4）微波辅助消解法}
微波辅助消解法是一种较新的试样分解技术，可用于有机、生物试样的氧化分解，也可用于难溶无机材料的分解。该方法是在常规分解方法的基础上，利用微波的热效应和非热效应来促进试样分解。

微波热效应是指物质直接吸收微波能量而升温；微波非热效应则与分子极化及其在高频交变电磁场中的取向和振荡有关。在这两种效应共同作用下，试样表层不断被激发并破裂，从而加快溶解或熔解过程。由于微波穿透能力强、传播速度快，因此能够同时作用于溶液或固体中的大量分子，使体系整体快速升温，显著提高加热效率。

微波消解一般在密闭容器中进行，可形成较高的温度和压力，不仅能提高分解效率，而且能减少溶剂用量，并降低易挥发组分的损失。

\subsubsection{消除干扰}

\subsubsection*{1. 干扰产生的原因}

实际分析对象通常较为复杂，除待测组分外，还含有其他组分，这些组分可能对分析测定产生干扰，因此在正式测定前常需要先消除干扰。

\subsubsection*{2. 消除干扰的方法}

\paragraph{（1）加入掩蔽剂}
通过化学反应将干扰组分转化为不干扰测定的形式。常见方法有：
\begin{itemize}
    \item 配位掩蔽法
    \item 沉淀掩蔽法
    \item 氧化还原掩蔽法
\end{itemize}

\paragraph{（2）将待测组分与干扰组分分离}
通过分离操作使待测组分与干扰组分彼此分开。其基本要求是：
\begin{itemize}
    \item 待测组分的损失应可忽略；
    \item 干扰组分应尽可能分离彻底。
\end{itemize}

常用分离方法有：
\begin{itemize}
    \item 沉淀分离法
    \item 溶剂萃取分离法
    \item 离子交换分离法
    \item 色谱分离法
\end{itemize}
\subsubsection{分析测定}
\subsubsection*{1. 分析方法的选择依据}

分析方法的选择应根据以下因素综合确定：
\begin{itemize}
    \item 待测组分的性质
    \item 待测组分的大致含量
    \item 对分析结果准确度的要求
    \item 方法本身的准确度、灵敏度、选择性和适用范围
\end{itemize}

\subsubsection*{2. 常见分析方法的特点}

\begin{itemize}
    \item \textbf{化学分析法：}准确度较高，适合常量分析。
    \item \textbf{仪器分析法：}灵敏度较高，适合微量组分分析。
\end{itemize}

\subsubsection*{3. 选择原则}

分析方法的选择一般应遵循以下原则：
\begin{itemize}
    \item 分析操作尽量简单、快速；
    \item 准确度应满足分析要求；
    \item 兼顾试剂与方法的经济性。
\end{itemize}
\subsubsection{分析结果的计算}
分析结果通常用待测组分的含量来表示。根据试样的取用量、测定所得数据以及分析过程中有关反应的计量关系，可以计算出试样中待测组分的含量。

\subsubsection*{1. 常用表示方法}

\begin{itemize}
    \item \textbf{质量分数：}待测组分的质量占试样质量的百分数，用于表示待测组分含量。
    \item \textbf{液体试样：}常用单位为 \(\mathrm{mg/L}\)。
    \item \textbf{气体试样：}常用单位为 \(\mathrm{mg/m^3}\)。
\end{itemize}
\newpage
\section{分析化学中的误差以及数据处理}
\subsection{分析化学中的误差与偏差}
\subsubsection{误差产生的原因}
{\bf{系统误差}}由于某种固定原因造成的误差，具有一定的规律性和可重复性。系统误差的来源包括：
\begin{itemize}
\item 仪器误差：如仪器的零点漂移、刻度不准确等。解决方案：定期校准仪器，使用标准物质进行校正。
\item 试剂误差：试剂蒸馏水的纯度不够，试剂中含有杂质等。解决方案：使用高纯试剂，定期检查试剂的质量。
\item 方法误差不恰当的实验设计或者分析方法的选择不当等。解决方案：选择合适的方法或者对方法矫正。
\item 操作误差：分析人员的操作和个人误差。
\item 单向性，重现性，可校正。由固定原因造成总是以相同的正负形式出现，大小可测量，能够校正
\end{itemize}

课件中还将系统误差细化为：
\begin{itemize}
    \item \textbf{试剂等级带来的差异：}优级纯（G.R.）、分析纯（A.R.）、化学纯（C.P.）、实验试剂（L.R.）及基准试剂。精密分析应优先选择高纯级别试剂与高纯水。
\end{itemize}

除系统误差外，还应区分：
\begin{itemize}
    \item \textbf{随机误差：}由电压、气流、温度、湿度等偶然因素引起，特点是不可校正、无法避免、服从统计规律。
    \item \textbf{过失（错误）：}由粗心引起（如称样洒落、记错读数），过失不能当作误差处理。
\end{itemize}
\subsubsection{误差的表示方法}
绝对误差：测量值与真值之差的绝对值。

相对误差：绝对误差与真值的比值，通常以百分数表示。

真值某一物体本身具有的客观存在的真实属性。但绝对真值不可测。（纯物质的理论值，标准物质证书上给的值，有经验的人用可靠方法多次测量取平均值确认消除系统误差（相对真值））
\subsubsection{描述实验结果}
实验结果的描述通常包括以下几个方面：
\begin{itemize}
\item 确定最有代表性的数据。
\item 描述一组数据的变化

\textbf{1. 平均偏差（Average Deviation）}

$$D = \frac{\sum_{i=1}^{n}|x_i - \bar{x}|}{n}$$

其中，$x_i$ 为各次测量值，$\bar{x}$ 为测量值的算术平均值，$n$ 为测量次数。

\textbf{2. 相对平均偏差（Relative Average Deviation）}

$$D_r = \frac{D}{\bar{x}} \times 100\%$$

相对平均偏差用于评估平均偏差相对于平均值的大小。

\textbf{3. 标准偏差（Standard Deviation）}

$$S = \sqrt{\frac{\sum_{i=1}^{n}(x_i - \bar{x})^2}{n-1}}$$

标准偏差反映了数据离散程度，更准确地描述了测量误差。$n-1$ 称为自由度。

\textbf{4. 相对标准偏差（Relative Standard Deviation，RSD）}

$$RSD= \frac{S}{\bar{x}} \times 100\%$$

相对标准偏差用于衡量标准偏差相对于平均值的百分比，也称为变异系数（CV）。

\textbf{5. 极差（全距）}

$$R=x_{\max}-x_{\min}$$

极差计算简单，但不能利用全部测量数据。
\end{itemize}
准确度与精确度：
\begin{itemize}
\item 准确度：测量值与真值的接近程度。准确度高表示测量值接近真值。
\item 精确度：测量值之间的一致程度。精确度高表示测量值之间的差异较小。
\end{itemize}
\subsection{提高结果准确度的方法}
\subsubsection{选择恰当的分析方法}
选择方法时应综合考虑：分析要求（准确度与灵敏度）、组分含量、实验室条件、方法的方便性与测试成本。

常见含量分级：常量组分（$>1\%$）、微量组分（$0.01\%\sim1\%$）、痕量组分（$<0.01\%$）。

\begin{itemize}
    \item 化学分析法：准确度较高（约 $0.1\%\sim0.2\%$），常用于常量组分。
    \item 仪器分析法：灵敏度高，常用于微量组分。
\end{itemize}
\subsubsection{减小测量误差——选择合适的测量对象的体积和质量}
分析天平读数误差常按 $\pm0.0001\,\mathrm{g}$ 估计；若采用差减称量，最大绝对误差可按 $\pm0.0002\,\mathrm{g}$ 估计。

课件给出：若要使称量相对误差小于 $0.1\%$，称样量应大于 $0.2\,\mathrm{g}$（最小称样量）。

滴定体积读数误差常按 $\pm0.01\,\mathrm{mL}$ 估计，差减读数最大绝对误差约 $\pm0.02\,\mathrm{mL}$。当体积相对误差要求小于 $0.1\%$ 时，体积宜大于 $20\,\mathrm{mL}$（常取 $20\sim30\,\mathrm{mL}$）。
\subsubsection{消除系统误差：对照试验、空白实验、仪器校准和分析结果的校正}
除对照试验、空白实验和仪器校准外，课件还强调\textbf{加标回收法}：
\[
\text{加标回收率}=\frac{c_2-c_1}{t}\times100\%
\]
其中 $c_1$ 为未加标样测定值，$c_2$ 为加标样测定值，$t$ 为理论加标量。

经验范围：常量组分回收率约 $99\%\sim101\%$，微量组分约 $90\%\sim110\%$。

注意事项：加标量应尽量接近样品待测物含量；加入体积一般小于 $1\%$ 以避免稀释效应；加标量通常不超过待测物含量的 $3$ 倍；加标物形态应与待测物形态相同。
\subsubsection{减小随机误差}多测几次
\subsection{误差的传递}

设分析结果为 \(R=f(x_1,x_2,\cdots,x_n)\)。

\subsubsection{系统误差的传递}
若各变量的系统误差分别为 \(\Delta x_1,\Delta x_2,\cdots,\Delta x_n\)，则结果的系统误差可按全微分近似表示为：

$$
\Delta R=\frac{\partial R}{\partial x_1}\Delta x_1+\frac{\partial R}{\partial x_2}\Delta x_2+\cdots+\frac{\partial R}{\partial x_n}\Delta x_n
$$

常用形式如下：
\begin{itemize}
    \item 若 \(R=a\pm b\)，则
$$
\Delta R=\Delta a+\Delta b
$$

    \item 若 \(R=ab\) 或 \(R=\dfrac{a}{b}\)，则
$$
\frac{\Delta R}{R}=\frac{\Delta a}{a}+\frac{\Delta b}{b}
$$

    \item 若 \(R=a^n\)，则
$$
\frac{\Delta R}{R}=n\frac{\Delta a}{a}
$$
\end{itemize}

\subsubsection{随机误差的传递}
若各变量的随机误差以标准偏差 \(s_{x_1},s_{x_2},\cdots,s_{x_n}\) 表示，且各变量相互独立，则结果的随机误差为：

$$
s_R=\sqrt{\left(\frac{\partial R}{\partial x_1}s_{x_1}\right)^2+\left(\frac{\partial R}{\partial x_2}s_{x_2}\right)^2+\cdots+\left(\frac{\partial R}{\partial x_n}s_{x_n}\right)^2}
$$

常用形式如下：
\begin{itemize}
    \item 若 \(R=a\pm b\)，则
$$
s_R=\sqrt{s_a^2+s_b^2}
$$

    \item 若 \(R=ab\) 或 \(R=\dfrac{a}{b}\)，则
$$
\frac{s_R}{R}=\sqrt{\left(\frac{s_a}{a}\right)^2+\left(\frac{s_b}{b}\right)^2}
$$

    \item 若 \(R=a^n\)，则
$$
\frac{s_R}{R}=|n|\frac{s_a}{a}
$$
\end{itemize}

\subsection{有效数字及其运算规则}

\subsubsection{有效数字的基本概念}

有效数字是指：
\begin{quote}
用来表示量的大小，同时反映测量准确程度的数字。
\end{quote}

在分析化学中，有效数字通常由以下两部分组成：
\[
\text{有效数字}=\text{全部可靠数字}+\text{一位不确定数字}
\]

\paragraph{例：}
\begin{itemize}
    \item 称量某样品 $1.2\,\mathrm{g}$，其不确定程度约为 $\pm 0.1\,\mathrm{g}$；
    \item 称量某样品 $1.2000\,\mathrm{g}$，其不确定程度约为 $\pm 0.0001\,\mathrm{g}$。
\end{itemize}

这说明：表示结果时，小数位越多，通常反映测量精度越高。

\subsubsection{确定有效数字的原则}

\begin{enumerate}
    \item 一个量值只保留 \textbf{一位不确定数字}。
    
    \item 从左边第一个不是 $0$ 的数字开始计数，之后的数字都属于有效数字，因此夹在中间的 $0$ 或末尾的 $0$ 都可能是有效数字。
    
    \paragraph{例：}
    \[
    0.0035 \text{（2位）},\qquad
    0.0530 \text{（3位）},\qquad
    0.3500 \text{（4位）},\qquad
    0.05030 \text{（4位）}
    \]
    
    \item 不能因为改变单位而改变有效数字的位数。
    
    \paragraph{例：}
    \[
    2.3\,\mathrm{L}=2.3\times 10^{3}\,\mathrm{mL}
    \]
    \[
    1.0\,\mathrm{mL}=0.0010\,\mathrm{L}
    \]
    
    \item 常数、系数、计数所得的自然数，其有效数字位数可认为是无限多位。
    
    \paragraph{例：}
    \[
    \pi,\quad e
    \]
    
    \item 对数、$\mathrm{pH}$ 的有效数字位数由 \textbf{小数部分}决定，整数部分只表示方次。
    
    \paragraph{例：}
    \[
    \mathrm{pH}=10.28 \qquad \text{（有效数字为 2 位）}
    \]
    \[
    [\mathrm{H}^+]=5.2\times 10^{-11}\,\mathrm{mol/L}
    \]
\end{enumerate}

\subsubsection{修约规则}

\paragraph{1. 四舍六入五成双}
修约时遵循以下规则：
\begin{itemize}
    \item 被修约数字小于 $5$ 时，舍去；
    \item 被修约数字大于 $5$ 时，进位；
    \item 被修约数字等于 $5$ 时，采用“五成双”原则：
    \begin{itemize}
        \item 若 $5$ 后面还有数字，则进位；
        \item 若 $5$ 后面没有数字，则看前一位：
        \begin{itemize}
            \item 前一位为奇数，进位；
            \item 前一位为偶数，舍去。
        \end{itemize}
    \end{itemize}
\end{itemize}

\paragraph{例：将下列数修约为 4 位有效数字}
\[
0.32474 \rightarrow 0.3247
\]
\[
0.32476 \rightarrow 0.3248
\]
\[
0.32475 \rightarrow 0.3248
\]
\[
0.32485 \rightarrow 0.3248
\]
\[
0.324851 \rightarrow 0.3249
\]

\paragraph{2. 禁止分次修约}
修约时应一步完成，不能先修约一次，再继续修约。

\paragraph{例：}
\[
0.5749 \not\rightarrow 0.575 \rightarrow 0.58
\]
正确应为：
\[
0.5749 \rightarrow 0.57
\]

\subsubsection{有效数字的运算规则}

\begin{enumerate}
    \item \textbf{加减法：}以小数点后位数最少的数据为准。
    
    \paragraph{例：}
    \[
    0.112+12.1+0.3214=12.5334
    \]
    按规则应保留到小数点后 1 位，因此结果写为：
    \[
    12.5
    \]
    
    \item \textbf{乘除法：}以有效数字位数最少的数据为准。
    
    \paragraph{例：}
    \[
    0.0121\times 25.66\times 1.0578=0.328432
    \]
    其中 $0.0121$ 有 3 位有效数字，因此结果应保留 3 位有效数字：
    \[
    0.328
    \]
\end{enumerate}

\paragraph{说明：}
用计算器计算时，中间过程可先不修约，最后结果再统一修约。

\subsubsection{分析化学中有效数字的书写要求}

\begin{enumerate}
    \item 组成含量大于或等于 $10\%$ 时，一般保留 4 位有效数字；
    \item 组成含量在 $1\%\sim 10\%$ 时，一般保留 3 位有效数字；
    \item 组成含量小于 $1\%$ 时，一般保留 2 位有效数字；
    \item 各种误差、偏差一般保留 $1\sim 2$ 位有效数字。
\end{enumerate}

\paragraph{课件中的示例：}
分析天平（称至 $0.1\,\mathrm{mg}$）：
\[
12.8228\,\mathrm{g}(6),\qquad 0.2348\,\mathrm{g}(4),\qquad 0.0600\,\mathrm{g}(3)
\]

滴定管（量至 $0.01\,\mathrm{mL}$）：
\[
26.32\,\mathrm{mL}(4),\qquad 3.97\,\mathrm{mL}(3)
\]

\subsubsection{本节笔记总结}

\begin{itemize}
    \item 有效数字既表示数值大小，也反映测量精度；
    \item 有效数字由“可靠数字 + 一位不确定数字”组成；
    \item 改变单位不能改变有效数字位数；
    \item 对数和 $\mathrm{pH}$ 的有效数字看小数部分；
    \item 修约采用“四舍六入五成双”；
    \item 修约必须一次完成，不能分次修约；
    \item 加减法看小数位数，乘除法看有效数字位数；
    \item 计算过程中中间值可多保留几位，最终结果再修约。
\end{itemize}


\subsection{分析化学中的数据处理}

数理统计是分析化学的重要分支，主要用于对实验数据进行整理、归纳和解释，为判断分析结果的可靠性提供依据。其基本思想是：在消除系统误差的前提下，用统计方法研究随机误差的分布规律及实验结果的波动特征。

\subsubsection{基本术语}

\begin{itemize}
    \item \textbf{总体（母体）}：被考察对象某一统计特征值的全体。
    \item \textbf{样本（子样）}：从总体中随机抽取的一组测量值。
    \item \textbf{样本容量}：样本中所含测量值的个数，记为 $n$。
\end{itemize}

\subsubsection{总体平均值与样本平均值}

设样本数据为 $x_1,x_2,\cdots,x_n$，则样本平均值为
\[
\bar{x}=\frac{x_1+x_2+\cdots+x_n}{n}=\frac{1}{n}\sum_{i=1}^{n}x_i
\]

当 $n\to\infty$ 时，总体平均值为
\[
\mu=\lim_{n\to\infty}\frac{1}{n}\sum_{i=1}^{n}x_i
\]

在确认系统误差已经消除的前提下，总体平均值可视为真值的估计。

\subsubsection{总体标准偏差与样本标准偏差}

样本标准偏差为
\[
s=\sqrt{\frac{\sum_{i=1}^{n}(x_i-\bar{x})^2}{n-1}}
\]

总体标准偏差为
\[
\sigma=\sqrt{\frac{\sum_{i=1}^{n}(x_i-\mu)^2}{n}}
\]

其中，$s$ 用于描述样本数据的离散程度，$\sigma$ 用于描述总体数据的离散程度。

\subsection{随机误差及其分布}

\subsubsection{随机误差的正态分布}

当测量次数足够多时，随机误差通常服从正态分布，其概率密度函数为
\[
y=f(x)=\frac{1}{\sigma\sqrt{2\pi}}
\exp\left[-\frac{(x-\mu)^2}{2\sigma^2}\right]
\]

其中：
\begin{itemize}
    \item $x$ 为测量值；
    \item $\mu$ 为总体平均值；
    \item $\sigma$ 为总体标准偏差；
    \item $y$ 为测量值出现的概率密度。
\end{itemize}

\subsubsection{正态分布的特点}

\begin{enumerate}
    \item 曲线关于 $\mu$ 对称，呈钟形。
    \item 随机误差绝对值越小，出现概率越大；绝对值越大，出现概率越小。
    \item 正误差与负误差出现的概率相等。
    \item $\sigma$ 越小，曲线越窄越高，说明数据越集中，分析结果精密度越高。
\end{enumerate}

\subsubsection{正态分布中的概率区间}

对于正态分布：
\[
P(|x-\mu|\le \sigma)\approx 68.3\%
\]
\[
P(|x-\mu|\le 2\sigma)\approx 95.5\%
\]
\[
P(|x-\mu|\le 3\sigma)\approx 99.7\%
\]

这表明：
\begin{itemize}
    \item 大误差出现的概率远小于小误差；
    \item 大多数测量值集中在平均值附近。
\end{itemize}

\subsubsection{少量实验数据的统计处理}

当测量次数较少（通常 $n<20$）时，随机误差不再直接按正态分布处理，而采用 \textbf{$t$ 分布} 进行统计分析。

\subsubsection{$t$ 分布的基本表达式}

样本平均值的标准偏差为
\[
s_{\bar{x}}=\frac{s}{\sqrt{n}}
\]

$t$ 值定义为
\[
t=\frac{\bar{x}-\mu}{s_{\bar{x}}}
=\frac{\bar{x}-\mu}{s/\sqrt{n}}
\]

自由度为
\[
f=n-1
\]

\subsubsection{$t$ 分布的特点}

\begin{itemize}
    \item $t$ 分布关于 $0$ 对称。
    \item 自由度越小，曲线越平缓、尾部越宽。
    \item 随着自由度增大，$t$ 分布逐渐接近正态分布。
    \item 当样本较少时，用 $t$ 分布比直接用正态分布更合理。
\end{itemize}

\subsection{平均值的置信区间}

在一定置信度下，总体平均值 $\mu$ 所在区间可表示为
\[
\mu=\bar{x}\pm t_{\alpha,f}\frac{s}{\sqrt{n}}
\]

其中：
\begin{itemize}
    \item $P$ 为置信度；
    \item $\alpha=1-P$ 为显著性水平；
    \item $f=n-1$ 为自由度；
    \item $t_{\alpha,f}$ 由 $t$ 分布表查得。
\end{itemize}

\subsubsection{置信区间的意义}

\begin{itemize}
    \item 在给定置信度 $P$ 下，总体平均值有较大概率落在该区间内。
    \item 置信度越高，区间通常越宽。
    \item 样本标准偏差越小、测量次数越多，置信区间越窄，结果越可靠。
\end{itemize}

\subsection{实验结果的比较与显著性检验}

显著性检验用于判断两个结果之间的差异是否仅由随机误差引起，还是已经达到“显著差异”的程度。

\subsubsection{显著性差异的含义}

\begin{itemize}
    \item \textbf{有显著性差异}：说明差异不能仅用随机误差解释，往往提示存在系统误差。
    \item \textbf{无显著性差异}：说明差异可归因于随机误差。
\end{itemize}

按检验方向可分为：
\begin{itemize}
    \item \textbf{双侧检验}：只关心是否存在差异，不区分正负方向。
    \item \textbf{单侧检验}：检验是否大于或小于某特定值，可分为右单侧与左单侧检验。
\end{itemize}

\subsubsection{$t$ 检验法：平均值与标准值的比较}

当需要比较测量平均值 $\bar{x}$ 与标准值 $\mu$ 是否一致时，采用
\[
t=\frac{|\bar{x}-\mu|}{s}\sqrt{n}
\]

判断方法如下：
\begin{enumerate}
    \item 计算实验 $t$ 值；
    \item 根据显著性水平 $\alpha$ 和自由度 $f=n-1$ 查表得临界值 $t_{\alpha,f}$；
    \item 若
    \[
    t>t_{\alpha,f}
    \]
    则说明存在显著性差异，可能有系统误差；
    \item 若
    \[
    t<t_{\alpha,f}
    \]
    则说明无显著性差异。
\end{enumerate}

分析化学中常以 $95\%$ 置信度作为检验标准，此时
\[
\alpha=0.05
\]

\subsubsection{$F$ 检验法：两组数据精密度的比较}

若要比较两组数据的精密度是否存在显著差异，可用 $F$ 检验。设两组数据标准偏差分别为 $s_1$、$s_2$，并约定 $s_1>s_2$，则
\[
F=\frac{s_1^2}{s_2^2}
\]

自由度分别为
\[
f_1=n_1-1,\qquad f_2=n_2-1
\]

判断方法：
\begin{enumerate}
    \item 计算 $F$ 值；
    \item 由显著性水平和自由度查 $F$ 表得临界值 $F_{\alpha}(f_1,f_2)$；
    \item 若
    \[
    F>F_{\alpha}(f_1,f_2)
    \]
    则说明两组数据精密度有显著差异；
    \item 若
    \[
    F<F_{\alpha}(f_1,f_2)
    \]
    则说明两组数据精密度无显著差异。
\end{enumerate}

\subsubsection{两组平均值的比较}

比较两组测定结果的平均值是否有显著差异时，通常先进行 $F$ 检验；若两组精密度无显著差异，再进行合并方差的 $t$ 检验。

合并标准偏差为
\[
S_A=
\sqrt{
\frac{(n_1-1)s_1^2+(n_2-1)s_2^2}{n_1+n_2-2}
}
\]

对应的统计量为
\[
t_A=
\frac{|\bar{x}_1-\bar{x}_2|}{S_A}
\sqrt{\frac{n_1n_2}{n_1+n_2}}
\]

自由度为
\[
f=n_1+n_2-2
\]

判断方法：
\begin{itemize}
    \item 若
    \[
    t_A>t_{\alpha,f}
    \]
    则两组平均值有显著差异；
    \item 若
    \[
    t_A<t_{\alpha,f}
    \]
    则两组平均值无显著差异。
\end{itemize}

\subsection{可疑数据的取舍}

实验中因偶然操作失误、仪器异常或读数错误，可能会出现明显偏离其余数据的可疑值。对可疑数据的取舍应采用统计方法判断，不能凭主观随意舍弃。

\subsubsection{$4\bar{d}$ 法}

平均偏差定义为
\[
\bar{d}=\frac{1}{n}\sum_{i=1}^{n}|x_i-\bar{x}|
\]

若某一测定值满足
\[
|x_i-\bar{x}|>4\bar{d}
\]
则该数据可视为可疑值，可考虑舍去。

\subsubsection{Grubbs 检验法}

Grubbs 检验适用于判断一组数据中某个极端值是否应舍去。

\subsubsection*{步骤}

\begin{enumerate}
    \item 将数据按大小规律排列：
    \[
    x_1,x_2,\cdots,x_n
    \]
    \item 若怀疑最大值 $x_n$ 为可疑值，则计算
    \[
    T=\frac{x_n-\bar{x}}{s}
    \]
    \item 若怀疑最小值 $x_1$ 为可疑值，则计算
    \[
    T=\frac{\bar{x}-x_1}{s}
    \]
    \item 查 Grubbs 临界值表，得 $T_{n,\alpha}$；
    \item 若
    \[
    T>T_{n,\alpha}
    \]
    则该值可舍去；否则应保留。
\end{enumerate}

\subsubsection{$Q$ 检验法}

$Q$ 检验法适用于较少数据中极端值的判断。

\subsubsection*{步骤}

\begin{enumerate}
    \item 将数据从小到大排列：
    \[
    x_1,x_2,\cdots,x_n
    \]
    \item 若怀疑最大值 $x_n$ 为可疑值，则
    \[
    Q=\frac{x_n-x_{n-1}}{x_n-x_1}
    \]
    \item 若怀疑最小值 $x_1$ 为可疑值，则
    \[
    Q=\frac{x_2-x_1}{x_n-x_1}
    \]
    \item 根据样本数和置信度查 $Q$ 表，得临界值 $Q_{n,P}$；
    \item 若
    \[
    Q>Q_{n,P}
    \]
    则该值可舍去；否则应保留。
\end{enumerate}

\subsubsection{可疑数据处理的注意事项}

\begin{itemize}
    \item 一般只对明显异常的单个数据进行判断。
    \item 不应连续随意舍去多个数据。
    \item 在正式统计检验前，应先处理可疑值。
\end{itemize}

\subsection{统计检验的一般顺序}

实验数据统计分析时，常按以下顺序进行：

\[
\text{可疑数据取舍} \;\longrightarrow\; F\text{检验} \;\longrightarrow\; t\text{检验}
\]

即：
\begin{enumerate}
    \item 先判断是否存在应舍去的可疑值；
    \item 再比较两组数据的精密度是否一致；
    \item 最后比较平均值之间是否存在显著差异。
\end{enumerate}

\subsubsection{本部分知识点总结}

\begin{enumerate}
    \item 基本统计量：总体、样本、样本容量、平均值、标准偏差。
    \item 大量数据的随机误差服从正态分布。
    \item 少量数据的统计处理采用 $t$ 分布。
    \item 平均值的可靠性可用置信区间表示。
    \item 显著性检验主要包括 $t$ 检验和 $F$ 检验。
    \item 可疑数据的取舍常用 $4\bar{d}$ 法、Grubbs 检验法和 $Q$ 检验法。
\end{enumerate}
\section{}
\end{document}
